% -*- coding: utf-8 -*-
% !TEX program = xelatex
\documentclass[plain,noanswer,medmath]{../../template/jnuexam}
\usepackage{../../template/kysx}

\begin{document}


\examtitle{2016}{2}


\loadproblems{1}{2016P1}%载入试卷一


\makepart{选择题}{1～8小题，每小题4分，共32分}


\begin{problem}
设$a_1 = x(\cos \sqrt{x} - 1)$，$a_2 = \sqrt{x} \ln(1 + \sqrt[3]{x})$，$a_3 = \sqrt[3]{x + 1} - 1$．
当$x \to 0^+$时，以上$3$个无穷小量按照从低阶到高阶的排序是\pickout{}
\begin{abcd}
\item $a_1,a_2,a_3$．
\item $a_2,a_3,a_1$．
\item $a_2,a_1,a_3$．
\item $a_3,a_2,a_1$．
\end{abcd}
\end{problem}


\useproblem{1}{1}{2}%【同试卷一第一(2)题】


\begin{problem}
反常积分①$\int_{-\infty}^0 \frac{1}{x^2}\e^{\frac{1}{x}} \dx$，
②$\int_0^{\text{+}\infty} \frac{1}{x^2}\e^{\frac{1}{x}} \dx$的敛散性为\pickout{}
\begin{abcd}
\item ①收敛，②收敛．
\item ①收敛，②发散．
\item ①发散，②收敛．
\item ①发散，②发散．
\end{abcd}
\end{problem}


\begin{problem}
设函数$f(x)$在$(-\infty , +\infty)$内连续，其导函数的图形如图所示，则\pickout{}
\par\centerline{\includegraphics[width=0.5\linewidth]{2016B}}
\begin{abcd}
\item 函数$f(x)$有$2$个极值点，曲线$y = f(x)$有$2$个拐点．
\item 函数$f(x)$有$2$个极值点，曲线$y = f(x)$有$3$个拐点．
\item 函数$f(x)$有$3$个极值点，曲线$y = f(x)$有$1$个拐点．
\item 函数$f(x)$有$3$个极值点，曲线$y = f(x)$有$2$个拐点．
\end{abcd}
\end{problem}


\begin{problem}
设函数$f_i(x)$（$i = 1,2$）具有二阶连续导数，且$f_i(x_0) < 0$（$i = 1,2$），
若两条曲线$y = f_i(x)$（$i = 1,2$）在点$(x_0,y_0)$处具有公切线$y = g(x)$，
且在该点处曲线$y = f_1(x)$的曲率大于曲线$y = f_2(x)$的曲率，则在$x_0$的某个邻域内，有\pickout{}
\begin{abcd}
\item $f_1(x) \le f_2(x) \le g(x)$．
\item $f_2(x) \le f_1(x) \le g(x)$．
\item $f_1(x) \le g(x) \le f_2(x)$．
\item $f_2(x) \le g(x) \le f_1(x)$．
\end{abcd}
\end{problem}


\begin{problem}
已知函数$f(x,y) = \frac{\e^x}{x - y}$，则\pickout{}
\begin{abcd}
\item $f'_x - f'_y = 0$．
\item $f'_x + f'_y = 0$．
\item $f'_x - f'_y = f$．
\item $f'_x + f'_y = f$．
\end{abcd}
\end{problem}


\useproblem{1}{1}{5}%【同试卷一第一(5)题】


\begin{problem}
设二次型$f(x_1,x_2,x_3) = a(x_1^2 + x_2^2 + x_3^2) + 2x_1x_2 + 2x_2x_3 + 2x_1x_3$
的正负惯性指数分别为$1,2$，则\pickout{}
\begin{abcd}
\item $a > 1$．
\item $a < - 2$．
\item $- 2 < a < 1$．
\item $a = 1$或$a = - 2$．
\end{abcd}
\end{problem}


\makepart{填空题}{9～14小题，每小题4分，共24分}


\begin{problem}
曲线$y = \frac{x^3}{1 + x^2} + \arctan(1 + x^2)$的斜渐近线方程为 \fillin{}．
\end{problem}


\begin{problem}
极限$\lim_{n \to \infty} \frac{1}{n^2}\left(\sin\frac{1}{n} + 2\sin\frac{2}{n}
+ \cdots + n\sin\frac{n}{n}\right) =$ \fillin{}．
\end{problem}


\begin{problem}
以$y = x^2 - \e^x$和$y = x^2$为特解的一阶非齐次线性微分方程为 \fillin{}．
\end{problem}


\begin{problem}
已知函数$f(x)$在$(-\infty,+\infty)$上连续，且$f(x) = (x + 1)^2 + 2\int_0^x f(t) \dt$，
则当$n \ge 2$时，$f^{(n)}(0) =$ \fillin{}．
\end{problem}


\begin{problem}
已知动点$P$在曲线$y = x^3$上运动，记坐标原点与点$P$间的距离为$l$．
若点$P$的横坐标对时间的变化率为常数$v_0$，则当点$P$运动到点$(1,1)$时，$l$对时间的变化率是 \fillin{}．
\end{problem}


\begin{problem}
设矩阵$\begin{pmatrix}
   a & -1 & -1 \\
  -1 &  a & -1 \\
  -1 & -1 & a
\end{pmatrix}$与$\begin{pmatrix}
  1 &  1 & 0 \\
  0 & -1 & 1 \\
  1 &  0 & 1
\end{pmatrix}$等价，则$a =$ \fillin{}．
\end{problem}


\makepart{解答题}{15～23小题，共94分}


\begin{problem}[points=10]
求极限$\lim_{x \to 0} (\cos 2x + 2x\sin x)^{\frac{1}{x^4}}$．
\end{problem}


\begin{problem}[points=10]
设函数$f(x) = \int_0^1 \left| t^2 - x^2 \right|\dt$（$x > 0$），求$f'(x)$并求$f(x)$的最小值．
\end{problem}


\begin{problem}[points=10]
已知函数$z = z(x,y)$由方程$(x^2 + y^2)z + \ln z + 2(x + y + 1) = 0$确定，求$z = z(x,y)$的极值．
\end{problem}


\begin{problem}[points=10]
设$D$是由直线$y = 1$，$y = x$，$y = - x$围成的有界区域，\goodbreak
计算二重积分$\iint_D \frac{x^2 - xy - y^2}{x^2 + y^2}\dx\dy.$
\end{problem}


\begin{problem}[points=10]
已知$y_1(x) = \e^x$，$y_2(x) = u(x) \e^x$是二阶微分方程$(2x - 1)y^n - (2x + 1)y' + 2y = 0$的解，
若$u(- 1) = e$，$u(0) = - 1$，求$u(x)$，并写出该微分方程的通解．
\end{problem}


\begin{problem}[points=11]
设$D$是由曲线$y = \sqrt{1 - x^2}$（$0 \le x \le 1$）与$\begin{cases}
  x = \cos^3 t \\
  y = \sin^3 t
\end{cases}$（$0 \le t \le \frac{\pi}{2}$）围成的平面区域，求$D$绕$x$轴旋转一周所得旋转体的体积和表面积．
\end{problem}


\begin{problem}[points=11]
已知$f(x)$在$\left[0,\frac{3\pi}{2} \right]$上连续，
在$\left(0,\frac{3\pi}{2} \right)$内是函数$\frac{\cos x}{2x - 3\pi}$的一个原函数$f(0) = 0$．
\begin{enumerate}
\item 求$f(x)$在区间$\left[0,\frac{3\pi}{2} \right]$上的平均值；
\item 证明$f(x)$在区间$\left(0,\frac{3\pi}{2} \right)$内存在唯一零点．
\end{enumerate}
\end{problem}


\begin{problem}[points=11]
设矩阵$A = \begin{pmatrix}
      1 & 1 & 1 - a \\
      1 & 0 & a \\
  a + 1 & 1 & a + 1
\end{pmatrix}$，$\beta = \begin{pmatrix}
  0 \\
  1 \\
  2a - 2
\end{pmatrix}$，且方程组$Ax = \beta$无解．\par
\begin{enumerate*}
\item 求$a$的值；
\item 求方程组$A^T Ax = A^T\beta$的通解．
\end{enumerate*}
\end{problem}


\useproblem[points=11]{1}{3}{21}%【同试卷一第三(21)题】


\end{document}
